\documentclass[UTF8,11pt]{ctexart}
\usepackage[a4paper,margin=24mm]{geometry}
\usepackage{amsmath,amssymb,amsthm,mathtools,mathrsfs}
\usepackage[all]{xy}
\usepackage{xurl,hyperref,enumitem}
\usepackage{tikz}
\usetikzlibrary{arrows.meta,cd,decorations.markings,calc,patterns}
\hypersetup{hidelinks,breaklinks=true}
\setlength{\emergencystretch}{3em}
\allowdisplaybreaks
\newtheorem*{theorem}{Theorem}
\newtheorem*{thm}{Theorem}
\newtheorem*{lemma}{Lemma}
\newtheorem*{proposition}{Proposition}
\newtheorem*{prop}{Proposition}
\newtheorem*{corollary}{Corollary}
\newtheorem*{cor}{Corollary}
\newtheorem*{claim}{Claim}
\theoremstyle{definition}
\newtheorem*{definition}{Definition}
\newtheorem*{defn}{Definition}
\newtheorem*{remark}{Remark}
\newtheorem*{example}{Example}
\newcommand{\R}{\mathbb R}
\newcommand{\C}{\mathbb C}
\newcommand{\Z}{\mathbb Z}
\newcommand{\Q}{\mathbb Q}
\newcommand{\N}{\mathbb N}
\newcommand{\F}{\mathbb F}
\newcommand{\D}{\mathbb D}
\newcommand{\bB}{\mathbb B}
\newcommand{\bC}{\mathbb C}
\newcommand{\bR}{\mathbb R}
\newcommand{\bZ}{\mathbb Z}
\newcommand{\bN}{\mathbb N}
\newcommand{\bQ}{\mathbb Q}
\newcommand{\bD}{\mathbb D}
\newcommand{\bF}{\mathbb F}
\newcommand{\CP}{\mathbb{CP}}
\newcommand{\RP}{\mathbb{RP}}
\newcommand{\sO}{\mathscr O}
\newcommand{\sI}{\mathscr I}
\newcommand{\sF}{\mathscr F}
\newcommand{\sG}{\mathscr G}
\newcommand{\sH}{\mathscr H}
\newcommand{\sR}{\mathscr R}
\newcommand{\sM}{\mathscr M}
\newcommand{\sV}{\mathscr V}
\newcommand{\sU}{\mathscr U}
\DeclareMathOperator{\Hom}{Hom}
\DeclareMathOperator{\Ext}{Ext}
\DeclareMathOperator{\Tor}{Tor}
\DeclareMathOperator{\Spec}{Spec}
\DeclareMathOperator{\Proj}{Proj}
\DeclareMathOperator{\supp}{supp}
\DeclareMathOperator{\im}{im}
\DeclareMathOperator{\id}{id}
\DeclareMathOperator{\Tr}{Tr}
\DeclareMathOperator{\tr}{tr}
\DeclareMathOperator{\rank}{rank}
\DeclareMathOperator{\ord}{ord}
\DeclareMathOperator{\dist}{dist}
\DeclareMathOperator{\End}{End}
\DeclareMathOperator{\Aut}{Aut}
\DeclareMathOperator{\Gal}{Gal}
\DeclareMathOperator{\vol}{vol}
\DeclareMathOperator{\Cl}{Cl}
\DeclareMathOperator{\coker}{coker}
\DeclareMathOperator{\codim}{codim}
\DeclareMathOperator{\divergence}{div}
\DeclareMathOperator{\Ric}{Ric}
\DeclareMathOperator{\grad}{grad}
\DeclareMathOperator{\Hess}{Hess}
\newcommand{\abs}[1]{\left\lvert#1\right\rvert}
\newcommand{\sabs}[1]{\lvert#1\rvert}
\newcommand{\babs}[1]{\bigl\lvert#1\bigr\rvert}
\newcommand{\Babs}[1]{\Bigl\lvert#1\Bigr\rvert}
\newcommand{\bbabs}[1]{\biggl\lvert#1\biggr\rvert}
\newcommand{\BBabs}[1]{\Biggl\lvert#1\Biggr\rvert}
\newcommand{\norm}[1]{\left\lVert#1\right\rVert}
\newcommand{\snorm}[1]{\lVert#1\rVert}
\newcommand{\bnorm}[1]{\bigl\lVert#1\bigr\rVert}
\newcommand{\Bnorm}[1]{\Bigl\lVert#1\Bigr\rVert}
\newcommand{\bbnorm}[1]{\biggl\lVert#1\biggr\rVert}
\newcommand{\BBnorm}[1]{\Biggl\lVert#1\Biggr\rVert}
\newcommand{\widebar}[1]{\overline{#1}}
\newcommand{\nicefrac}[2]{\frac{#1}{#2}}
\newcommand{\linnprod}[2]{\langle#1,#2\rangle}
\newcommand{\myindex}[1]{#1}
\newcommand{\glsadd}[1]{}
\newcommand{\avoidbreak}{}
\newcommand{\sourceref}[1]{\textnormal{[原文：\nolinkurl{#1}]}}
\newcommand{\thmref}[1]{\sourceref{#1}}
\newcommand{\lemmaref}[1]{\sourceref{#1}}
\newcommand{\propref}[1]{\sourceref{#1}}
\newcommand{\corref}[1]{\sourceref{#1}}
\newcommand{\defnref}[1]{\sourceref{#1}}
\newcommand{\exampleref}[1]{\sourceref{#1}}
\newcommand{\exerciseref}[1]{\sourceref{#1}}
\newcommand{\figureref}[1]{\sourceref{#1}}
\newcommand{\sectionref}[1]{\sourceref{#1}}
\newcommand{\appendixref}[1]{\sourceref{#1}}
\newcommand{\stacksref}[2]{\href{https://stacks.math.columbia.edu/tag/#1}{#2 [#1]}}


\newcommand{\E}{\mathbb E}
\newcommand{\Prob}{\mathbb P}
\newcommand{\Reff}{\mathcal R_{\mathrm{eff}}}
\newcommand{\eps}{\varepsilon}
\DeclareMathOperator{\diam}{diam}
\DeclareMathOperator{\Var}{Var}
\DeclareMathOperator{\Crit}{Crit}
\newcommand{\dd}{\,\mathrm d}

\begin{document}
\section*{065\quad 平坦模的Lazard有限自由模滤过余极限表示}
\subsection*{来源与计数目标}
The Stacks Project Authors, \emph{The Stacks Project}, Section 10.81. 主命题 Theorem 10.81.4 (058G)，及 Lemmas 10.81.1--10.81.2。实际访问版本：2026-09-28。
\par\url{https://stacks.math.columbia.edu/tag/058C}
\par\textbf{本文件只计一个问题：}任意 $R$-模平坦，当且仅当它是有限自由 $R$-模有向系统的余极限。
\subsection*{作者命题与人类证明}
\begin{lemma}[10.81.1]
Let $M$ be an $R$-module. The following are equivalent:
\begin{enumerate}
\item $M$ is flat.
\item If $f:R^n\to M$ is a module map and $x\in\operatorname{Ker}(f)$, then there are module maps $h:R^n\to R^m$ and $g:R^m\to M$ such that $f=g\circ h$ and $x\in\operatorname{Ker}(h)$.
\item Suppose $f:R^n\to M$ is a module map, $N\subset\operatorname{Ker}(f)$ any submodule, and $h:R^n\to R^m$ a map such that $N\subset\operatorname{Ker}(h)$ and $f$ factors through $h$. Then given any $x\in\operatorname{Ker}(f)$ we can find a map $h':R^n\to R^{m'}$ such that $N+Rx\subset\operatorname{Ker}(h')$ and $f$ factors through $h'$.
\item If $f:R^n\to M$ is a module map and $N\subset\operatorname{Ker}(f)$ is a finitely generated submodule, then there are module maps $h:R^n\to R^m$ and $g:R^m\to M$ such that $f=g\circ h$ and $N\subset\operatorname{Ker}(h)$.
\end{enumerate}
\end{lemma}
\begin{proof}
That (1) is equivalent to (2) is just a reformulation of the equational criterion for flatness. To show (2) implies (3), let $g:R^m\to M$ be the map such that $f$ factors as $f=g\circ h$. By (2) find $h'':R^m\to R^{m'}$ such that $h''$ kills $h(x)$ and $g:R^m\to M$ factors through $h''$. Then taking $h'=h''\circ h$ works. (3) implies (4) by induction on the number of generators of $N\subset\operatorname{Ker}(f)$ in (4). Clearly (4) implies (2).
\end{proof}

\begin{lemma}[10.81.2]
Let $M$ be an $R$-module. Then $M$ is flat if and only if the following condition holds: if $P$ is a finitely presented $R$-module and $f:P\to M$ a module map, then there is a free finite $R$-module $F$ and module maps $h:P\to F$ and $g:F\to M$ such that $f=g\circ h$.
\end{lemma}
\begin{proof}
This is just a reformulation of condition (4) from Lemma 10.81.1.
\end{proof}

\begin{theorem}[10.81.4, Lazard's theorem, Tag 058G]
Let $M$ be an $R$-module. Then $M$ is flat if and only if it is the colimit of a directed system of free finite $R$-modules.
\end{theorem}
\begin{proof}
A colimit of a directed system of flat modules is flat, as taking directed colimits is exact and commutes with tensor product. Hence if $M$ is the colimit of a directed system of free finite modules then $M$ is flat.

For the converse, first recall that any module $M$ can be written as the colimit of a directed system of finitely presented modules, in the following way. Choose a surjection $f:R^I\to M$ for some set $I$, and let $K$ be the kernel. Let $E$ be the set of ordered pairs $(J,N)$ where $J$ is a finite subset of $I$ and $N$ is a finitely generated submodule of $R^J\cap K$.

Then $E$ is made into a directed partially ordered set by defining $(J,N)\leq(J',N')$ if and only if $J\subset J'$ and $N\subset N'$. Define $M_e=R^J/N$ for $e=(J,N)$, and define $f_{ee'}:M_e\to M_{e'}$ to be the natural map for $e\leq e'$. Then $(M_e,f_{ee'})$ is a directed system and the natural maps $f_e:M_e\to M$ induce an isomorphism $\operatorname{colim}_{e\in E}M_e\xrightarrow{\cong}M$.

Now suppose $M$ is flat. Let $I=M\times\mathbf Z$, write $(x_i)$ for the canonical basis of $R^I$, and take in the above discussion $f:R^I\to M$ to be the map sending $x_i$ to the projection of $i$ onto $M$. To prove the theorem it suffices to show that the $e\in E$ such that $M_e$ is free form a cofinal subset of $E$. So let $e=(J,N)\in E$ be arbitrary.

By Lemma 10.81.2 there is a free finite module $F$ and maps $h:R^J/N\to F$ and $g:F\to M$ such that the natural map $f_e:R^J/N\to M$ factors as $R^J/N\xrightarrow{h}F\xrightarrow{g}M$. We are going to realize $F$ as $M_{e'}$ for some $e'\geq e$.

Let $\{b_1,\ldots,b_n\}$ be a finite basis of $F$. Choose $n$ distinct elements $i_1,\ldots,i_n\in I$ such that $i_\ell\notin J$ for all $\ell$, and such that the image of $x_{i_\ell}$ under $f:R^I\to M$ equals the image of $b_\ell$ under $g:F\to M$. This is possible since every element of $M$ can be written as $f(x_i)$ for infinitely many distinct $i\in I$ (by our choice of $I$).

Now let $J'=J\cup\{i_1,\ldots,i_n\}$, and define $R^{J'}\to F$ by $x_i\mapsto h(x_i)$ for $i\in J$ and $x_{i_\ell}\mapsto b_\ell$ for $\ell=1,\ldots,n$. Let $N'=\operatorname{Ker}(R^{J'}\to F)$. Observe:
\begin{enumerate}
\item The square
\[\xymatrix{R^{J'}\ar[r]\ar@{^{(}->}[d]&F\ar[d]^g\\R^I\ar[r]_f&M}\]
is commutative, hence $N'\subset K=\operatorname{Ker}(f)$;
\item $R^{J'}\to F$ is a surjection onto a free finite module, hence it splits and so $N'$ is finitely generated;
\item $J\subset J'$ and $N\subset N'$.
\end{enumerate}
By (1) and (2) $e'=(J',N')$ is in $E$, by (3) $e'\geq e$, and by construction $M_{e'}=R^{J'}/N'\cong F$ is free.
\end{proof}
\subsection*{整理说明（非作者正文）}
主证明的两个方向均取入，特别保留 $I=M\times\mathbf Z$ 的重复生成元构造、$J',N'$ 的实现及共尾性论证；不是只给出范畴表达。正文中的 $R^I$ 是原文用来表示以 $I$ 为基的自由模的记号，按作者记号保留；并非在此新增直积约定。

标准先修为平坦性的等式判据（10.39.11）、有向余极限的正合性、张量积与有向余极限交换、自由模的投射性。取材不含无关的 10.81.3；10.81.1--2 仅为本题支撑，不重复计数。

难度依据：要把消去有限关系的局部条件组织为有限自由对象组成的整个有向系统，并真正实现一个共尾子系统；仅知道平坦性的定义不能直接得到此表示。外部引用保留原编号，正文来自网页数学 TeX，另存转录副本而非原章节下载件。
\subsection*{可修改的简短标注（非作者正文）}
$c$：一般模的平坦性与表示。

$a$：平坦性的等式判据；有限呈示；自由模分解；关系消去；有向系统；共尾子集；张量积与余极限。

\texttt{cross\_theory\_status = yes}。把平坦性的关系消去条件转成经有限自由模的因子分解，再通过新生成元构造共尾有向系统，得到模的全局余极限表示。位置：10.81.1--2 与 10.81.4 的逆向构造。
标签为非穷尽、可修改初稿。
\subsection*{许可与修改记录}
原作者及作品见来源栏；来源采用 GNU Free Documentation License 1.2 或后续版本。本题证摘录和附加整理沿用该许可。2026-09-28：增加题号、来源定位、独立排版、整理说明和初稿标注；数学内容的取材范围及排版变化见整理说明。
\end{document}
